% In compendium/laws-of-harmony-and-dynamics/law-of-harmonic-shifting.tex
\subsection*{Law of Harmonic Shifting}
\hrule
\vspace{0.3cm}
\GetLaw{LawOfHarmonicShifting}
\vspace{0.5cm}

\subsubsection*{Conceptual Overview}
This principle describes how stable systems react to external influence.  A perturbation does not destroy the system's stability; rather, it causes the system to gracefully shift its entire orbital configuration to a new, different, but equally stable state. 

\subsubsection*{Proof}
The proof involves applying a transformation operator to a stable system and demonstrating that the resulting system satisfies the Law of Algebraic Stability under a new set of coordinates. 

\subsubsection*{Formal Proof}
The proof relies on the property of linearity. We will show that if a linear operator $L$ is applied to all points in a stable system, the resulting system's Center of Gravity remains zero.

\begin{enumerate}
    \item Let a system be defined by a set of points $\{p_i\}$ and corresponding harmonic weights $\{w_i\}$. The system is stable, so according to the \textbf{Law of Algebraic Stability}, its Center of Gravity is zero:
    $$ \Xi_{\text{initial}} = \sum_{i} w_i p_i = 0 $$
    
    \item Let $L$ be an arbitrary linear transformation operator. We apply this operator to every point $p_i$ in the system to create a new set of shifted points, $\{p'_i\}$, where $p'_i = L(p_i)$.
    
    \item We now calculate the Center of Gravity of this new, shifted system, $\Xi_{\text{shifted}}$:
    $$ \Xi_{\text{shifted}} = \sum_{i} w_i p'_i = \sum_{i} w_i L(p_i) $$
    
    \item Because the operator $L$ is linear, it satisfies the distributive property, so we can factor it out of the summation:
    $$ \Xi_{\text{shifted}} = L\left(\sum_{i} w_i p_i\right) $$
    
    \item From step 1, we know that the expression inside the parentheses is the Center of Gravity of the initial system, which is zero. Therefore:
    $$ \Xi_{\text{shifted}} = L(0) = 0 $$
    
    \item The Center of Gravity of the shifted system is zero. Thus, the new system is also stable according to the \textbf{Law of Algebraic Stability}. The law is hereby \textbf{proven}.
\end{enumerate}

\subsubsection*{Computational Verification}
The following Python script symbolically verifies this principle. It constructs an explicitly stable system, applies a generic linear operator, and confirms that the resulting system's Center of Gravity is also zero.

\begin{lstlisting}[language=Python, caption={Symbolic proof of the Law of Harmonic Shifting.}, label={lst:harmonic_shifting_proof}]
import sympy


def prove_harmonic_shifting():
    """
    Symbolically proves the Law of Harmonic Shifting.

    This script demonstrates that if an arbitrary linear operator is applied
    to every point in an algebraically stable system, the resulting
    (shifted) system is also algebraically stable.
    """
    # --- Define symbolic variables ---
    # T, J are harmonic weights. p1 is an arbitrary initial point.
    T, J, p1 = sympy.symbols('T J p1')

    # 'c' represents an arbitrary linear transformation operator (e.g., a complex scalar)
    c = sympy.symbols('c')

    print("=" * 60)
    print("Symbolic Proof of the Law of Harmonic Shifting")
    print("=" * 60)

    # --- Step 1: Construct an initially stable system ---
    # From the Law of Algebraic Stability, T*p1 + J*p2 = 0.
    # We solve for p2 to guarantee this condition.
    p2 = -(T / J) * p1

    # The Center of Gravity of the initial system is T*p1 + J*p2.
    # By construction, this is zero.
    initial_cg = sympy.simplify(T * p1 + J * p2)

    print("1. An initially stable system (p1, p2) is constructed.")
    print(f"   p1 = {p1}")
    print(f"   p2 = {p2}")
    print(f"   Initial Center of Gravity (T*p1 + J*p2) simplifies to: {initial_cg}\n")
    assert initial_cg == 0

    # --- Step 2: Apply a linear transformation 'c' to each point ---
    p1_shifted = c * p1
    p2_shifted = c * p2

    print("2. A generic linear operator 'c' is applied to each point.")
    print(f"   p1_shifted = {p1_shifted}")
    print(f"   p2_shifted = {p2_shifted}\n")

    # --- Step 3: Calculate the Center of Gravity of the new, shifted system ---
    shifted_cg = T * p1_shifted + J * p2_shifted

    print("3. Calculating the Center of Gravity of the new (shifted) system...")
    print(f"   New C.G. = T * (p1_shifted) + J * (p2_shifted)")
    print(f"   New C.G. = {shifted_cg}\n")

    # --- Step 4: Verify that the new system is also stable ---
    simplified_shifted_cg = sympy.simplify(shifted_cg)

    print("4. Verifying the stability of the new system.")
    print(f"   The new Center of Gravity simplifies to: {simplified_shifted_cg}\n")
    assert simplified_shifted_cg == 0

    print("Conclusion: The Center of Gravity for the shifted system is 0.")
    print("The Law of Harmonic Shifting is computationally verified.")
    print("=" * 60)


if __name__ == '__main__':
    prove_harmonic_shifting()
\end{lstlisting}

\paragraph{Verification Output}
\begin{verbatim}
============================================================
Symbolic Proof of the Law of Harmonic Shifting
============================================================
1. An initially stable system (p1, p2) is constructed.
    p1 = p1
    p2 = -T*p1/J
    Initial Center of Gravity (T*p1 + J*p2) simplifies to: 0

2. A generic linear operator 'c' is applied to each point.
    p1_shifted = c*p1
    p2_shifted = -T*c*p1/J

3. Calculating the Center of Gravity of the new (shifted) system...
    New C.G. = T * (p1_shifted) + J * (p2_shifted)
    New C.G. = 0

4. Verifying the stability of the new system.
    The new Center of Gravity simplifies to: 0

Conclusion: The Center of Gravity for the shifted system is 0.
The Law of Harmonic Shifting is computationally verified.
============================================================
\end{verbatim}